\documentclass[11pt,letterpaper]{article}
\usepackage[T1]{fontenc}
\usepackage[utf8]{inputenc}
\usepackage{lmodern}
\usepackage[margin=0.65in]{geometry}
\usepackage{xcolor}
\usepackage{hyperref}
\usepackage{fancyhdr}
\renewcommand{\familydefault}{\sfdefault}
\definecolor{ink}{HTML}{202A26}
\definecolor{muted}{HTML}{505A55}
\hypersetup{colorlinks=true,urlcolor=ink,pdftitle={Wanglong Lu - ML Engineer Resume},pdfauthor={Wanglong Lu}}
\ifdefined\pdfgentounicode\input{glyphtounicode}\pdfgentounicode=1\fi
\setlength{\parindent}{0pt}
\setlength{\parskip}{1.5pt}
\setlength{\headheight}{14pt}
\pagestyle{fancy}
\fancyhf{}
\renewcommand{\headrulewidth}{0pt}
\fancyfoot[L]{\footnotesize\color{muted}Wanglong Lu}
\fancyfoot[R]{\footnotesize\color{muted}\thepage\ / 2}
\newcommand{\cvsection}[1]{\par\addvspace{6pt}{\large\bfseries #1}\par\vspace{1pt}\hrule\vspace{3pt}}
\newcommand{\company}[2]{\par\addvspace{4pt}\textbf{#1}\hfill{\small\color{muted}#2}\par}
\newcommand{\position}[2]{\par\addvspace{2pt}\textbf{#1}\hfill{\small #2}\par}
\newenvironment{points}{\begin{list}{\textbullet}{\setlength{\leftmargin}{12pt}\setlength{\itemsep}{1.2pt}\setlength{\parsep}{0pt}\setlength{\topsep}{1.5pt}\setlength{\partopsep}{0pt}}}{\end{list}}
\newcommand{\project}[2]{\par\addvspace{3pt}\textbf{#1}\hfill{\small #2}\par}
\begin{document}
\fontsize{10.5}{12.7}\selectfont
\raggedright
{\fontsize{25}{28}\selectfont\bfseries Wanglong Lu}\par
{\large Senior Data Scientist \textbar\ ML Systems \& AI Engineering}\par
St. John's, NL, Canada \textbar\ \href{mailto:lwlxhl@gmail.com}{lwlxhl@gmail.com} \textbar\
\href{https://longlongaaago.github.io/}{Website} \textbar\
\href{https://github.com/LonglongaaaGo}{GitHub} \textbar\
\href{https://www.linkedin.com/in/wanglong-lu-8a98b6235/}{LinkedIn} \textbar\
\href{https://scholar.google.com/citations?user=TuxCf4UAAAAJ}{Google Scholar}

\cvsection{Profile}
Computer Science Ph.D.; \textbf{30 scholarly works overall} (author-reported). Builds \textbf{AWS-backed ML tools and AI applications}, including \textbf{SageMaker job submission} and an \textbf{internally used RAG assistant}; develops models in PyTorch.

\cvsection{Technical Skills}
\textbf{Engineering:} Python, SQL, Git, Bash, C++, Docker; FastAPI, Mangum, PostgreSQL.\par
\textbf{Cloud \& Compute:} AWS SageMaker, Lambda, Bedrock, S3; PyTorch, CUDA, SLURM.\par
\textbf{ML \& GenAI:} Feature engineering, Transformers, diffusion, PEFT/LoRA, Hugging Face, RAG, agentic tools.

\cvsection{Professional Experience}
\company{Nasdaq}{St. John's, NL, Canada}
\position{Senior Data Scientist, AI/Analytics}{Mar. 2025 -- Present}
\begin{points}
\item Delivered \textbf{MLP application integration} into a \textbf{SageMaker workspace} from an existing cross-team UI/backend prototype; configured \textbf{networking and IAM}. Platform at \textbf{internal deployment stage}.
\item Implemented \textbf{PostgreSQL-backed user-state persistence}; deployed \textbf{FastAPI services on AWS Lambda using Mangum} and connected them to the platform application layer.
\item Built \textbf{SageMaker Job Submitter for MLP}: script/dependency packaging for \textbf{Processing, Training, HPO and Pipelines}; cut \textbf{new-task configuration/submission preparation from 7 days to 1 day (about 86\%)}.
\item Developed \textbf{PagerDuty RAG assistant v2} for incident-context retrieval and troubleshooting guidance; \textbf{now used internally and recognized in a company shout-out}.
\item Designed and implemented \textbf{feature engineering with PyTorch/CUDA-accelerated feature-label AUROC screening}; used \textbf{Amazon Bedrock} to assess candidates' semantic plausibility.
\item Designed cross-scale ensemble feature selection for \textbf{cheque-fraud ML (CFML)}, removing \textbf{77.2\% of input features} while retaining competitive performance in evaluations.
\item Mentored projects in \textbf{input payload checking, GNN-based crime-ring detection, cheque OCR and tabular-data generation}; the latter led to a \href{https://doi.org/10.21203/rs.3.rs-7624993/v1}{\textbf{StDDPM manuscript submission}}.
\end{points}

\position{AI Algorithm Intern (two placements)}{May 2024 -- Jan. 2025}
\begin{points}
\item Designed and implemented \href{https://arxiv.org/abs/2503.04021}{\textbf{TextDoctor}} in \textbf{Python/PyTorch} for document/cheque inpainting; ran AWS/SLURM/CUDA experiments and demonstrated \textbf{8K inpainting on a 24GB RTX 3090}.
\item Evaluated \textbf{7 document datasets} without dataset-specific fine-tuning. FUNSD PaddleOCR word accuracy: \textbf{55.38\% vs. DocDiff's 45.24\% (+10.14 percentage points)} using TextDoctor (GSDM).
\item \textbf{First-author TextDoctor manuscript submitted to EAAI}; the publicly linked version is an \textbf{arXiv preprint}, not an accepted journal article.
\item Presented research after \textbf{both internship placements}: about \textbf{40--60 attendees} at the earlier talk and \textbf{39 cross-functional attendees} at the final talk.
\end{points}

\company{Beaufort Solutions}{St. John's, NL, Canada}
\position{AI Intern}{Jul. 2022 -- Nov. 2022}
\begin{points}
\item Led design of \textbf{TEG}, a \textbf{CLIP-based few-shot adaptation} method for photo-book curation, using a text-embedding-guided classifier and auxiliary loss; evaluated with Python/PyTorch and AWS/SLURM.
\item Built the collection pipeline and led \textbf{Theme25 dataset construction: 35,655 images, 25 themes}; evaluated three datasets; research published in \href{https://doi.org/10.1117/1.JEI.33.1.013028}{\textbf{Journal of Electronic Imaging (2024)}}.
\end{points}

\company{Zhongkong Hanlian}{Hangzhou, China}
\position{AI Algorithm Intern}{Dec. 2017 -- Aug. 2018}
\begin{points}
\item \textbf{Deployed vehicle-logo recognition with Caffe at community security gates}; designed \textbf{ResNet--DenseNet fusion}, built Java annotation tools, and constructed a training/evaluation dataset.
\end{points}

\newpage
{\large\bfseries Wanglong Lu}\hfill{\small AI Tooling, Model Development \& Education}\par

\cvsection{Agentic AI Tooling}
\project{Local Agentic Code Editor}{Nasdaq internal exploratory prototype; 2026}
\begin{points}
\item Built a working local code editor connecting \textbf{language-model APIs} to \textbf{task-dependent routing across code generation, code review and shell execution}.
\item Implemented an iterative debugging loop: \textbf{diagnose bugs, edit code, apply fixes and inspect execution feedback} to determine the next action. Remains an exploratory prototype, not an approved production product.
\end{points}

\cvsection{Selected Model Development \& Publications}
\project{\href{https://doi.org/10.1109/TNNLS.2026.3707463}{UltraDiffEdit -- High-Resolution Generative Editing} {\small\normalfont[\href{https://github.com/LonglongaaaGo/UltraDiffEdit}{Code}]}}{First author; IEEE TNNLS 2026}
\begin{points}
\item Extended pretrained diffusion models to \textbf{tuning-free 8K editing on one RTX 3090} with multi-patch encoding and global-local sampling; tested SDXL/ControlNet on \textbf{three benchmarks (1,188 images)}.
\item \textbf{UHRSDEdit (4K--8K) crop-FID 7.45 vs. 18.76 for SDXL+bicubic (60.3\% lower)}.
\end{points}

\project{\href{https://doi.org/10.52202/085713-0664}{AdaMSS -- Parameter-Efficient Adaptation} {\small\normalfont[\href{https://github.com/jzheng20/AdaMSS}{Code}]}}{Second author; NeurIPS 2025}
\begin{points}
\item Co-developed multi-subspace tuning; \textbf{ViT-Large accuracy 90.13\% vs. LoRA's 85.40\%} across \textbf{7 tasks} with \textbf{15.4\% of LoRA's trainable parameters}. Method \href{https://github.com/huggingface/peft/tree/main/examples/adamss_finetuning}{\textbf{integrated into Hugging Face PEFT}}.
\end{points}

\project{\href{https://doi.org/10.1109/TCSVT.2026.3720433}{DocPure -- Unified Document Restoration} {\small\normalfont[\href{https://github.com/LingmingSSS/DocPure}{Code}]}}{Co-first author; IEEE TCSVT 2026}
\begin{points}
\item Co-developed prompt-free, structure-guided wavelet restoration for \textbf{4 tasks and 12 datasets}. On TDD deblurring vs. DocRes: \textbf{PSNR 32.55 vs. 28.34 dB; PaddleOCR word accuracy 61.27\% vs. 55.66\%}.
\end{points}

\project{\href{https://doi.org/10.1016/j.patcog.2024.111312}{Visual Style Prompt Learning} {\small\normalfont[\href{https://github.com/LonglongaaaGo/VSPBFR}{Code}]}}{First author; Pattern Recognition 2025}
\begin{points}
\item Developed diffusion style prompts: \textbf{four denoising steps, 12.3 FPS at 512\,$\times$\,512 on RTX 3080 Ti}; \textbf{WebPhoto-Test FID 74.25 vs. CodeFormer's 83.17 (10.7\% lower)}.
\item On \textbf{CelebA-Test}, FAN landmark NME \textbf{6.08\% to 2.43\%}; HSEmotion agreement \textbf{79.73\% to 86.03\%}, using the respective models' predictions on clean images as references.
\end{points}

\project{\href{https://doi.org/10.1109/TVCG.2024.3434386}{FACEMUG -- Multimodal Local Facial Editing}}{First author; TVCG 2025 (online 2024)}
\begin{points}
\item Designed feature fusion and latent warping for \textbf{six input modalities}; \textbf{77.88\% preference over Unite\&Conquer} in a \textbf{52-participant CelebA-HQ study} with sketch/semantic/text guidance.
\end{points}

\project{\href{https://doi.org/10.26599/CVM.2025.9450408}{GRIG -- Data-Efficient Image Inpainting} {\small\normalfont[\href{https://github.com/LonglongaaaGo/GRIG_few_shot_inpainting}{Code}]}}{First author; CVM 2025}
\begin{points}
\item Author-manuscript evaluation of CNN--Transformer residual inpainting: \textbf{10 datasets (18--750 training images)}; \textbf{3.70--26.48\% lower FID vs. next-best baselines at 50--60\% masking}.
\end{points}
{\small Full publication record: \href{https://scholar.google.com/citations?user=TuxCf4UAAAAJ}{Google Scholar} and \href{https://longlongaaago.github.io/publications/}{publication portfolio}.}

\cvsection{Education}
\textbf{Ph.D., Computer Science} -- Memorial University of Newfoundland, Canada \hfill 2021--2025\par
\textbf{M.Eng., Computer Software and Theory} -- Wenzhou University, China \hfill 2018--2021\par
\textbf{B.Eng., Digital Media Technology} -- Communication University of Zhejiang, China \hfill 2014--2018

\cvsection{Selected Recognition \& Service}
Outstanding Research Award (Ph.D., Memorial, 2023). \textbf{Five granted Chinese invention patents as co-inventor}; assignee at grant: Wenzhou University.\par
Mentored \textbf{12 students} in generative vision and applied ML, including restoration and recognition projects.\par
Teaching assistant in Data Structures and Algorithms (2021) and Introduction to Machine Learning (2024).\par
{\raggedright Invited reviewer for \textbf{IEEE TPAMI, TIP, TMM, TCSVT, JBHI and SPL}; \textbf{Pattern Recognition}, EAAI, JVCI, Knowledge-Based Systems, Neurocomputing, Applied Soft Computing, Scientific Reports, Displays; ECCV 2026.\par}
Presented feature-selection research to \textbf{60+ attendees} (2026); invited speaker, IEEE Newfoundland and Labrador Section AGM (2026); guest lecturer on deep generative models, Memorial University (2025).
\end{document}
